\documentclass[10pt]{article}

\usepackage[a4paper,margin=25mm]{geometry}
\usepackage{amsmath}
\usepackage{amssymb}
\usepackage{booktabs}
\usepackage{hyperref}
\usepackage{microtype}
\usepackage{xcolor}

\hypersetup{colorlinks=true,linkcolor=blue!50!black,citecolor=blue!50!black,urlcolor=blue!60!black}
\newcommand{\project}{Touhou PC-98 RL}
\newcommand{\todo}[1]{\textcolor{red!70!black}{[TODO: #1]}}

\title{CPU-First Reinforcement Learning for Native PC-98 Bullet-Hell Games\\
\large A Reproducible Systems Study and Research Plan}
\author{Touhou PC-98 RL Project}
\date{Working draft --- 22 August 2026}

\begin{document}
\maketitle

\begin{abstract}
We study reinforcement learning (RL) directly against a native PC-98 bullet-hell
game under a CPU-only constraint.  The principal systems problem is not merely
policy optimization: a useful experiment must boot a legally supplied disk
image, identify the correct emulated-memory region, provide isolated resets,
extract observations at game rate, and avoid CPU oversubscription.  We present
an end-to-end TH05 environment that uses exact child-process attachment, compact
entity-set observations, a 180,136-parameter recurrent actor--critic, and
multi-process proximal policy optimization (PPO).  Compact reads have a measured
median latency of 0.516 ms, compared with 3.550 ms for the dense-map path, while
reducing transferred floats by a factor of 776.  Eight emulators deliver about
223 transitions/s on the tested host.  After 16,384 training transitions and
validation-based checkpoint selection, an untouched four-seed, 900-step paired
test improves mean scalar return from 0.0529 to 0.3639, halves observed deaths
from four to two, and reduces mean resource cost by 17.4\%.  No evaluated run
clears the stage; these results show preliminary short-horizon improvement, not
convergence.  We document negative results, validity threats, and a path from
the current TH05 adapter toward a common TH01--05 research framework.
\end{abstract}

\section{Research questions}

This project is organized around five falsifiable research questions.

\begin{description}
  \item[RQ1: Systems feasibility.] Can an unmodified emulator plus structured
  memory extraction sustain useful RL throughput on CPU without interfering
  with other workloads on a shared host?

  \item[RQ2: Representation efficiency.] Can a compact, permutation-invariant
  representation of nearby hazards replace large, sparse spatial maps while
  preserving enough state for bullet-hell control?

  \item[RQ3: Learning signal.] Under a small CPU budget, does recurrent PPO
  improve survival and resource management over a paired untrained policy?

  \item[RQ4: Safe exploration.] Do adapter-certified hard action constraints
  improve sample efficiency without merely hiding an unsafe policy?

  \item[RQ5: Generality.] Which interfaces can be shared by Touhou 1--5, and
  which memory, termination, reward, and patching components must remain
  game-specific?
\end{description}

The present evidence addresses RQ1 strongly, RQ2 at the systems level, and RQ3
only preliminarily.  RQ4 is now implemented as an experimental ablation but has
no comparative result yet.  RQ5 is deliberately deferred until the agent can
complete TH05 on its highest difficulty; only TH05 is currently validated.

\section{Background and related work}

The initial reverse-engineering reference describes a high-fidelity,
multi-objective bullet-hell environment and an accelerator-oriented baseline
\cite{liu2026thrl}.  We treat that work as a source for low-level TH05 state
discovery, not as the architecture or identity of the present framework.

Our optimizer follows clipped PPO \cite{schulman2017ppo} with generalized
advantage estimation (GAE) \cite{schulman2015gae}.  A game frame contains a
variable-size collection of bullets and special projectiles.  Processing these
objects as sets follows the invariance motivation of Deep Sets
\cite{zaheer2017deepsets}; learned attention pooling is related to Set
Transformer \cite{lee2019settransformer}, although our encoder is deliberately
smaller and does not perform quadratic self-attention among all hazards.

DOSBox-X supplies PC-98 emulation \cite{dosboxx2026}.  ReC98 release P0335 is
used only as a clean boot baseline \cite{rec98p0335}; it is not a runtime or
build dependency.

\section{Failure analysis of the reference baseline}

The reference project solved valuable reverse-engineering problems, but its
released implementation encodes assumptions that do not hold for our CPU-only,
multi-process, current-emulator setting.  We audited those assumptions and
reproduced their consequences before replacing them.  Table~\ref{tab:failures}
separates observed failure modes from design preferences; ``inadequate'' here
means inadequate for the stated experimental setting, not universally wrong.

\begin{table}[ht]
  \centering
  \small
  \caption{Reference-baseline assumptions, observed consequences, and our
  replacements.}
  \label{tab:failures}
  \begin{tabular}{p{0.21\linewidth}p{0.32\linewidth}p{0.35\linewidth}}
    \toprule
    Baseline assumption & Why it fails here & Replacement \\
    \midrule
    XPU/CUDA-only package and dense CNN input & No supported GPU is available;
    each 2,048-step map rollout was documented as 1.62 GiB & CPU PyTorch lockfile,
    compact entity sets, and a small GRU \\
    At most 8--10 MiB per scanned memory region & The current DOSBox-X guest
    mapping is about 30 MiB, so game signatures are never searched & Writable
    anonymous-region scan with a tested 64 MiB safety cap \\
    Shared host game directory & Concurrent emulators can read and write the
    same configuration and score files & One private HDI copy and exact child
    PID per environment \\
    Dense $24\times92\times96$ maps every step & Allocation and language-boundary
    transfer dominate CPU inference & 273-float native compact reader \\
    Normalize every objective before weighting & Reward magnitude loses its
    semantic meaning and sparse/noisy objectives receive unit variance & Vector
    GAE, fixed semantic scalarization, then one normalization \\
    Overwrite a latest-only checkpoint & A later update can silently replace a
    better earlier policy & Periodic snapshots, validation LCB, untouched test \\
    No deadline for a live-but-stalled worker & The learner can wait forever
    while leaving emulator children behind & Bounded startup/step IPC and exact
    cleanup \\
    Readable memory implies rollout readiness & About 80 opening action slots
    ignore control but still emit survival reward & Reset gate waits for the
    player input lock to end \\
    Patch stage values are human numbered & The loader writes them directly to
    TH05's zero-based resident stage, silently shifting scenarios by one &
    Human Stage 1--6 mapped to resident 0--5; both values recorded \\
    \bottomrule
  \end{tabular}
\end{table}

\subsection{Version-fragile systems assumptions}

The released memory walker discards mappings larger than 10 MiB during its
initial search and uses an 8 MiB limit for subsequent readable-region filters.
On the tested DOSBox-X release, the relevant writable guest mapping is roughly
30 MiB.  This is a hard compatibility failure rather than a throughput issue:
the correct signature cannot be found in bytes that are never searched.  The
new 64 MiB cap is still bounded, but its value is checked against the actual
mapping layout.  Searches are further restricted to writable anonymous
regions, retaining a small candidate set.

The original spawn path correctly attaches to its newly created child, but all
workers mount one host \texttt{export} directory.  Its separate attach-existing
path chooses the first process returned by a name search.  The former risks
cross-worker file-state interference and the latter is ambiguous whenever
unrelated emulators run on the host.  Private writable disk images plus exact
child ownership make environment identity explicit.  This matters for science
as well as reliability: a reset is not reproducible if another worker can alter
its persistent state.

The same audit found unbounded coordination waits for workers that remain alive
but stop responding, and setup instructions that rely on privileged execution
and world-writable build directories.  Our runner instead uses child-process
tracing, finite deadlines, deterministic cleanup, and ordinary user
permissions.  These changes reduce both stalled experiments and accidental
interference with unrelated jobs.

Memory readability is also weaker than actionability.  In a live directional
probe, the initial player state was readable immediately but remained fixed for
about 80 action intervals while the opening invincibility/input lock was active.
The baseline collector records this prefix and attributes its survival reward
to sampled actions that the game cannot execute.  Our reset gate waits for both
an initialized player and the end of this lock.  A sustained-left/right probe
then moved normalized X from 0.50 to 0 and back beyond 0.62, independently
confirming that the subsequent action path is effective.

\subsection{CPU cost is determined by representation}

The reference package pins an XPU build of PyTorch and states that a GPU is
required.  More importantly, its learning path stores a
$24\times92\times96$ tensor for every transition, even though most cells are
empty.  Merely selecting a CPU wheel would therefore leave the principal
allocation, copy, memory-bandwidth, and CNN costs intact.  Our change is
algorithmic as well as infrastructural: nearby hazards are variable-size sets,
so a set encoder expresses the symmetry directly and makes CPU inference
smaller.  Table~\ref{tab:reader} reports the resulting measured reader cost.

\subsection{Optimization semantics and model selection}

For objective $o$, the reference trainer first constructs
$\widehat{A}^{(o)}=\operatorname{standardize}(A^{(o)})$, combines these
standardized values, and standardizes the sum again.  Consequently, multiplying
an objective's raw reward by almost any positive constant has no effect, and a
rare or noisy objective can be promoted to the same empirical variance as a
dense, reliable one.  This is not intrinsically invalid, but it no longer
implements scalarization in declared reward units.  We retain a vector critic
and vector GAE, apply explicit semantic weights to the unnormalized advantages,
and normalize the resulting policy advantage exactly once.

Finally, overwriting a single ``latest'' file assumes that training progress is
monotone.  Our short runs directly contradict that assumption: a policy that
improved at 24,576 transitions regressed badly by 49,152.  Periodic immutable
snapshots, a disjoint validation split, and an untouched paired test do not
eliminate selection bias, but they make it visible and prevent the last update
from becoming the result by accident.

\section{System design}

\subsection{Reproducible native-game execution}

The environment starts from a user-provided Japanese HDI.  A fail-closed patch
script verifies complete SHA-256 hashes before applying the required executable
layouts.  The supplied archive is a mixed revision: \texttt{MAIN.EXE} requires
patch layout 1, whereas \texttt{OP.EXE} matches layout 2.  Treating layout
numbers as a paired version produces a black screen.  Patched binaries and game
data are never committed.

Each Gymnasium environment copies the 21 MiB template HDI into a private
temporary directory.  This makes resets independent and prevents concurrent
workers from racing on game-written score or configuration files.  DOSBox-X is
spawned as a direct child, and the native extension attaches to that exact PID.
The memory search is restricted to writable anonymous mappings no larger than
64 MiB.  This includes the approximately 30 MiB guest-memory mapping used by the
tested DOSBox-X release; the previous 8--10 MiB cap silently excluded it.

Startup and step operations have explicit timeouts.  Reset retries cover both
process creation and the interval before initialized gameplay becomes readable.
Workers release inputs and terminate their exact child process on shutdown.

\subsection{Compact observation}

The compact observation has 273 floats:
\begin{equation}
  37\ \text{global features}
  + 16\!\times\!7\ \text{special-projectile features}
  + 16\!\times\!7\ \text{bullet features}
  + 12\ \text{item features}.
\end{equation}
Entities are ordered by distance for truncation, but the neural encoder pools
them as a set.  Each entity contains relative position, velocity, type, speed,
and normalized distance.  Missing entities use the sentinel
$(0,0,0,0,0,0,1)$.  A previously plausible mask based only on nonzero values is
wrong because the sentinel distance is one; the implementation now recognizes
the full sentinel and has a regression test for all-padding inputs.

The older observation additionally builds 24 maps of size $92\times96$, or
211,968 floats per step.  Our hot path constructs no map.  This is important
because map allocation and Python transfer can cost more than policy inference
on CPU.

\subsection{Policy and value model}

Special projectiles and bullets pass through separate shared token networks.
Each set is summarized by learned attention pooling and masked max pooling.
Global state and item features have small multilayer perceptrons (MLPs); the
four summaries are fused into a 128-dimensional vector.  A one-layer GRU with
128 hidden units models temporal context.  The actor predicts 19 discrete
actions.  A shared critic head predicts three value components rather than
collapsing objectives before value estimation.  The model contains 180,136
parameters.

\subsection{CPU recurrent PPO}

For scaled vector reward $\mathbf{r}_t\in\mathbb{R}^3$, vector GAE is
\begin{align}
  \boldsymbol{\delta}_t &= \mathbf{r}_t
  + \gamma(1-d_t)\mathbf{V}(s_{t+1}) - \mathbf{V}(s_t),\\
  \mathbf{A}_t &= \boldsymbol{\delta}_t
  + \gamma\lambda(1-d_t)\mathbf{A}_{t+1}.
\end{align}
The policy advantage is $A_t=\mathbf{w}^{\top}\mathbf{A}_t$, normalized once
after scalarization.  This avoids independently normalizing sparse objectives
and then assigning them equal empirical variance.  PPO uses ratio clipping,
value clipping, entropy regularization, gradient clipping, and approximate-KL
early stopping.

Rollouts are synchronous across independent emulator processes.  Every worker
uses one PyTorch thread; the learner defaults to eight.  The parent process is
CPU-affined and run with positive niceness.  Periodic snapshots are selected on
a validation seed set using
\begin{equation}
  S = \bar{R} - z\,\frac{s_R}{\sqrt{n}},
\end{equation}
after prioritizing observed stage successes.  Thus the last PPO update is not
automatically considered the best model.

Two post-baseline extensions are exposed for controlled ablation.  Analytic
geometry appends bounded constant-velocity closest-approach features using
physical scales supplied by the game adapter.  Hard safety masks only actions
the adapter can certify as unavailable or behaviorally redundant---currently a
bomb with zero stock and movement into an active boundary clamp.  The masked
categorical is renormalized identically during rollout and PPO optimization,
and removed probability mass is logged.  It does not implement scripted
collision avoidance.

\section{Experimental protocol}

\paragraph{Host.}
The test machine exposes 96 logical CPU threads and approximately 755 GiB RAM.
Other experiments occupied selected high-numbered cores.  RL jobs were run at
niceness 10 and confined to known-idle low-numbered cores.  No GPU was visible
to PyTorch.  The software stack used Python 3.14.6, PyTorch 2.13.0+cpu, and
DOSBox-X tag \texttt{dosbox-x-v2026.07.02}.

\paragraph{Game condition.}
The first short-run dataset actually starts at Stage 2 (resident stage index 1)
with three lives, three bombs, character 2, rank 0, and power 0.  It was
initially labelled Stage 1 by following the patch README; source audit showed
that the loader assigns \texttt{skip\_to} directly to TH05's zero-based resident
field.  We retain and relabel those measurements rather than discard or rewrite
them.  New scenario tooling accepts human Stage 1--6 and writes resident 0--5.
The environment advances at a 36 ms action interval.  Evaluations are stochastic
and paired by PyTorch sampling seed.  The game itself is not yet randomized
through a controlled seed API.

\paragraph{Training.}
The selected policy comes from training seed 202 at update 8: 16,384 total
transitions from eight workers.  Six snapshots were compared for 900 steps on
validation seeds 41, 43, and 44.  The final comparison uses untouched seeds
51--54.  These test seeds were not used for training or model selection.

\paragraph{Metrics.}
We report observation latency, transferred floats, model parameters, rollout
throughput, scalar return, death count, and the raw resource-reward component.
Stage success is the primary task metric but is absent from all short runs.
The exact aggregate and per-seed values are tracked in
\path{experiments/2026-08-22-th05-cpu.json}.

\section{Results}

\subsection{Systems efficiency}

\begin{table}[ht]
  \centering
  \caption{Paused-game observation-reader benchmark.}
  \label{tab:reader}
  \begin{tabular}{lrrr}
    \toprule
    Reader & Median (ms) & 95th percentile (ms) & Floats returned \\
    \midrule
    Dense maps & 3.550 & 5.014 & 211,968 + 273 \\
    Compact & 0.516 & 0.542 & 273 \\
    \bottomrule
  \end{tabular}
\end{table}

Compact reading is 6.88 times faster by median and transfers 776 times fewer
map-equivalent floats.  The compact actor--critic has 180,136 parameters versus
2,375,478 for the comparison CNN--GRU model, a 13.2-fold reduction.

Eight live environments sustain approximately 223 transitions/s when no reset
occurs.  A 2,048-transition rollout takes about 9.19 s, while three PPO epochs
take 0.20--0.33 s.  Therefore environment time, not gradient computation, is
currently dominant.  Increasing learner threads from one to eight improved the
model microbenchmark, but 16 threads caused severe oversubscription; bounded
threading is part of the algorithmic system, not a cosmetic setting.

\subsection{Snapshot selection}

\begin{table}[ht]
  \centering
  \caption{Validation results for the six candidate snapshots (three seeds,
  900 steps each).  LCB uses $z=1$.}
  \label{tab:selection}
  \begin{tabular}{lrrr}
    \toprule
    Candidate & Mean return & Standard error & LCB score \\
    \midrule
    Reproduction update 8 & 0.2333 & 0.2341 & -0.0007 \\
    Reproduction update 10 & 0.2112 & 0.2573 & -0.0461 \\
    Reproduction update 12 & 0.0167 & 0.0145 & 0.0022 \\
    Seed 101 update 8 & 0.2408 & 0.2379 & 0.0030 \\
    Seed 202 update 8 & 0.2743 & 0.2362 & \textbf{0.0381} \\
    Seed 303 update 8 & 0.2769 & 0.2401 & 0.0368 \\
    \bottomrule
  \end{tabular}
\end{table}

The candidate with the highest mean is not selected: seed 202 has a slightly
higher lower-confidence score than seed 303.  This illustrates why latest-only
or single-seed checkpointing is unreliable at this budget.

\subsection{Untouched paired test}

\begin{table}[ht]
  \centering
  \caption{Untouched stochastic test.  Resource values closer to zero are
  better.}
  \label{tab:test}
  \begin{tabular}{rrrrrrr}
    \toprule
    Seed & \multicolumn{3}{c}{Untrained} & \multicolumn{3}{c}{Selected PPO} \\
    \cmidrule(lr){2-4}\cmidrule(lr){5-7}
      & Return & Deaths & Resource & Return & Deaths & Resource \\
    \midrule
    51 & 0.0085 & 1 & -157 & 0.0135 & 1 & -157 \\
    52 & 0.0165 & 1 & -155 & 0.6875 & 0 & -87 \\
    53 & 0.1985 & 1 & -83  & 0.7325 & 0 & -69 \\
    54 & -0.0120 & 1 & -168 & 0.0220 & 1 & -152 \\
    \midrule
    Mean/total & 0.0529 & 4 & -140.75 & \textbf{0.3639} & \textbf{2} & \textbf{-116.25} \\
    \bottomrule
  \end{tabular}
\end{table}

The paired mean return gain is 0.3110 with standard error 0.1707.  Deaths fall
by 50\%, and mean resource cost improves by 17.4\%.  With only four seeds, the
uncertainty is large; a normal 95\% interval for return gain would cross zero.
The correct conclusion is preliminary improvement, not statistical certainty.

\subsection{Negative and unstable results}

A hand-written short-horizon collision heuristic performs worse than random in
the tested condition, reaching four deaths and terminal failure around
948--1,107 steps.  Independent nearest-bullet prediction does not capture
multi-projectile escape geometry; the heuristic is retained only as a negative
ablation and is not used for behavior cloning.

An earlier exploratory policy improved four-seed return after 24,576
transitions, then regressed badly by 49,152 transitions.  Three additional
16,384-transition seeds varied widely.  These observations motivated periodic
snapshots, validation-based selection, and an untouched test split.

\subsection{Matched Lunatic ablation and curriculum probe}

Plain compact recurrent PPO and the analytic-geometry plus hard-safety variant
were each trained with seeds 1011, 2022, and 3033 for 24,576 transitions.  We
retained updates 4, 8, and 12, selected nine candidates per condition on seeds
71--73, and tested the selected pair on untouched seeds 81--84 for 1,200 steps.
Source audit relabelled this condition as Stage 2 (resident index 1).

\begin{table}[ht]
  \centering
  \caption{Matched TH05 Lunatic Stage 2 ablation.  LCB is validation mean minus
  one standard error; deaths are totals across four test seeds.}
  \label{tab:lunatic-ablation}
  \begin{tabular}{lrrrr}
    \toprule
    Condition & Validation mean & LCB & Test mean & Test deaths \\
    \midrule
    Plain PPO & -0.856 & -1.449 & -1.648 & 14 \\
    Geometry + hard safety & -0.822 & -1.422 & \textbf{-1.054} & \textbf{11} \\
    \bottomrule
  \end{tabular}
\end{table}

Neither condition clears the stage.  The paired test gain is 0.594 with
standard error 0.669, and one seed contributes most of it.  Thus the point
estimate favors geometry plus safety but does not establish a reliable gameplay
improvement.  The shield constrains 40.1\% of test decisions while removing
only 2.47\% of policy probability on average, consistent with a narrow legality
constraint rather than a scripted dodge policy.

A true Stage 1 Lunatic probe at phase 3, ending at phase 4, is cleared by a
hard-safety random policy in 397 steps with one death.  This confirms that phase
termination provides a dense success signal.  It motivates a completion-gated
curriculum: learn no-miss completion on a short phase, move the starting phase
backward, reduce privileged power/resources, and only then optimize a full
stage.  The machine-readable record is
\path{experiments/2026-08-22-th05-lunatic-ablation.json}.

\section{Threats to validity}

\begin{itemize}
  \item Both paired tests have only four stochastic seeds (900 or 1,200 steps
  per seed).  Paired standard errors are large and no full-stage clear occurs.
  \item Evaluation seeds control action sampling, but the native game does not
  yet expose a separately controlled environment RNG seed.
  \item The comparison model has no reproducible compatible checkpoint, so
  systems comparisons do not imply a gameplay win over prior work.
  \item Learning evidence comes from one TH05 Stage 2 configuration; the Stage
  1 result is only a random-policy phase probe.  Neither establishes transfer
  to other stages, characters, ranks, or TH01--04.
  \item The compact observation drops most hazards beyond the nearest 16 in
  each category and currently omits a compact normal-enemy set.
  \item Host timing is measured while unrelated jobs are present.  CPU affinity
  and niceness reduce interference but do not produce a dedicated-machine
  benchmark.
\end{itemize}

\section{Research roadmap}

The next experiments are specified as hypotheses rather than assumed upgrades.

\paragraph{H1: richer compact geometry.}
Evaluate the implemented closest-approach features, then add normal-enemy
tokens while keeping the observation below 2 KiB.  Compare each addition
against the 273-float schema with matched training transitions and
validation/test splits.

\paragraph{H2: TH05 Lunatic mastery.}
Train a stage/phase curriculum that ends at rank 3 with no privileged starting
resources.  Promote a policy only after stage completion, then optimize no-miss
completion and finally full-game completion.  Use completion and
resource-normalized completion as primary metrics rather than short-prefix
return.

\paragraph{H3: reset-aware collection.}
When terminals become common, replace global synchronous waiting with an
asynchronous collector or reset pool.  The prediction is improved wall-clock
throughput without changing the on-policy age bound.

\paragraph{H4: risk-aware objectives.}
Compare the fixed scalarization with constrained or distributional survival
objectives.  Any method must report the full reward vector and may not hide a
survival regression behind score gain.

\paragraph{H5: adapter-certified safety.}
Compare no mask, exact resource/boundary constraints, and any future predictive
shield at matched training budgets.  Report removed policy probability mass
and evaluate every trained policy both with and without its shield.

\paragraph{H6: old-five-games interface.}
Define per-game contracts for observation schema, action meaning, terminal
flags, executable identity, and reset.  TH01--04 enter the benchmark only after
live parser/action tests pass and TH05 mastery has been addressed; shared code
alone is not evidence of support.

\section{Conclusion}

Native PC-98 RL is practical on a CPU when observation construction, process
attachment, worker isolation, and thread scheduling are treated as first-class
research concerns.  The present system removes a large dense-map bottleneck,
runs eight independent games at useful speed, and exhibits a promising but
small-data improvement in survival and resource use.  The main scientific task
now is not to declare success, but to turn this systems baseline into a larger,
pre-registered multi-configuration study with true stage completion metrics.

\bibliographystyle{plain}
\bibliography{references}

\end{document}
